\chapter*{Abstract}

This work explores the development and application of low-temperature scanning probe microscopy using nitrogen-vacancy (NV) centers in diamond as quantum sensors to investigate the magnetic properties of two-dimensional (2D) superconductors. The NV center serves as an optically addressable solid-state qubit, capable of high-sensitivity magnetometry with nanoscale spatial resolution. 

The experimental framework involves the construction of a Scanning Quantum Microscope (SQM) integrated with a closed-cycle cryogenic system. Operating NV centers at low temperatures introduces several experimental challenges that must be addressed. While numerous advanced measurement techniques are reported for NV centers, their integration into an SQM remains complex. This work details the development of a cryogenic SQM that overcomes these challenges, enabling robust operation. To enhance measurement performance, advanced modalities such as gradiometry, two-point tracking, and sequential Bayesian experiments is implemented.

A primary focus of this research is the investigation of vortex dynamics in few-atomic-layer 2D superconductors. Detecting vortices in these materials poses a challenge for conventional magnetic imaging techniques due to weak magnetic signals and the frequent surface contamination of exfoliated samples. Consequently, there remains a lack of comprehensive reports on the static and dynamic nature of vortices in these systems, despite the potential for magnetic imaging to significantly advance our understanding of 2D superconductivity. The developed SQM enables the direct visualization of the Meissner effect and the mapping of Pearl vortices, which exhibit broader magnetic profiles in thin films compared to bulk materials. Systematic studies revealed that vortex arrangements in these materials are highly dependent on the cooling rate, reflecting a competition between thermal fluctuations and disorder pinning as the system transitions from a vortex liquid to a vortex glass phase. Furthermore, investigations into step junctions demonstrate that these features act as barriers for vortices, leading to merged vortex lines and accumulation along the edges. Signatures of non-integer flux quantization are also observed.

The capabilities of the SQM are fundamentally driven by the availability and performance of its sophisticated scanning probes. These probes are built around diamond waveguide structures that optimize photon collection efficiency while delivering atomic force microscopy resolution. This work details advancements in diamond nanofabrication aimed at improving both the sensitivity and production yield of these scanning probes. A robust, high-yield fabrication process is developed utilizing PMMA, titanium hard masks, and inductively coupled reactive ion etching to create high-aspect-ratio diamond waveguides. Additionally, a novel island-waveguide design is introduced. This architecture localizes an ensemble of preferentially aligned NV centers strictly at the probe's apex, significantly enhancing photon collection and magnetic sensitivity while maintaining the spatial resolution necessary for precision scanning applications.

This work presents the development and improvement of an SQM probe, the development of a complete cryogenic SQM system, and the investigation of novel properties of a 2D superconductor using the same.